\chapter{Materialized Inference}
\label{ch:inference}

An ontology says more than it stores. If \code{Alloy subClassOf Metal} and
\code{Metal subClassOf Material}, then every alloy is a material, though no
triple says so; \cref{ch:engine} ended on exactly that gap, with a query for
materials that missed \code{Alloy}. A store that answers ontology queries has
to close it. This chapter treats inference as a database problem: the
entailed triples are a \emph{view} of the stored ones, and the questions are
the old ones for views: compute it on read, or store it and keep it current on
every write.

\section{Two ways to answer with entailments}

\begin{description}[leftmargin=1.2em]
\item[Rewrite at read time] (backward chaining). Turn \code{?x a :Material}
into \code{?x a ?c . ?c rdfs:subClassOf* :Material}, a property path, and let
the executor follow it. Nothing extra is stored; every query pays for the
closure it needs, and a query planner must estimate a recursive path.
\item[Materialize at write time] (forward chaining). Compute every entailed
triple once, store it in the same indexes, and answer \code{?x a :Material}
with one scan of POS, as if the triples had been asserted. Reads stay as fast
as they are; writes and space pay.
\end{description}

\begin{center}
\begin{tabular}{@{}lll@{}}
\toprule
& Rewrite & Materialize \\
\midrule
Read cost & a recursive path per query & one scan \\
Write cost & none & maintain the closure \\
Space & asserted triples & asserted plus entailed \\
Planner & estimate recursion & unchanged \\
\bottomrule
\end{tabular}
\end{center}

Ontologies are read far more often than they change, and their queries ask
for classes \emph{and} their subclasses almost every time, so OntoDB
materializes, like RDFox and GraphDB.

\section{The rules}

RDFS entailment is a handful of rules \cite{rdfsemantics}. Written as Datalog,
with \code{T(s,p,o)} for ``the store holds triple $(s,p,o)$'':
\begin{align*}
T(c, \mathit{sc}, e) &\leftarrow T(c, \mathit{sc}, d),\ T(d, \mathit{sc}, e) && \text{(rdfs11, subclass is transitive)}\\
T(x, \mathit{type}, d) &\leftarrow T(x, \mathit{type}, c),\ T(c, \mathit{sc}, d) && \text{(rdfs9, types go up)}\\
T(p, \mathit{sp}, r) &\leftarrow T(p, \mathit{sp}, q),\ T(q, \mathit{sp}, r) && \text{(rdfs5, subproperty is transitive)}\\
T(x, q, y) &\leftarrow T(x, p, y),\ T(p, \mathit{sp}, q) && \text{(rdfs7, triples go up)}\\
T(x, \mathit{type}, c) &\leftarrow T(x, p, y),\ T(p, \mathit{domain}, c) && \text{(rdfs2)}\\
T(y, \mathit{type}, c) &\leftarrow T(x, p, y),\ T(p, \mathit{range}, c) && \text{(rdfs3)}
\end{align*}
Each rule body is a two-pattern basic graph pattern, so each rule is one join
of two scans, and the executor of \cref{ch:joins} already evaluates it. The
closure is the least set of triples containing the data and closed under the
rules.

\section{Naive and semi-naive evaluation}

The obvious way to compute the closure is to apply every rule to everything,
add what is new, and repeat until a round adds nothing. That is
\term{naive evaluation}, and it is wasteful: round $i$ rederives every triple of
rounds $1$ to $i-1$ before finding the new ones.

The fix is to notice that a triple can be new in round $i$ only if one of the
triples it was derived from was new in round $i-1$. \term{Semi-naive
evaluation} \cite{alicebook} keeps the set $\Delta$ of triples new in the last
round and evaluates each rule with at least one body pattern restricted to
$\Delta$:
\[
\Delta_{i+1} \;=\; \Big(\textstyle\bigcup_{\text{rules}}\ \big(\Delta_i \Join T_i\big) \cup \big(T_i \Join \Delta_i\big)\Big) \setminus T_i,
\qquad T_{i+1} = T_i \cup \Delta_{i+1}.
\]
For a chain of subclasses $A \sqsubseteq B \sqsubseteq C \sqsubseteq D \sqsubseteq E$:

\begin{center}
\begin{tabular}{@{}lll@{}}
\toprule
Round & $\Delta$ (new subclass triples) & Joins evaluated \\
\midrule
0 & the 4 asserted links & --- \\
1 & $A{\sqsubseteq}C$, $B{\sqsubseteq}D$, $C{\sqsubseteq}E$ & $\Delta_0 \times T_0$ \\
2 & $A{\sqsubseteq}D$, $A{\sqsubseteq}E$, $B{\sqsubseteq}E$ & $\Delta_1 \times T_1$ \\
3 & (none) & $\Delta_2 \times T_2$; stop \\
\bottomrule
\end{tabular}
\end{center}

Naive evaluation would rejoin all four asserted links in every round. On a
deep hierarchy the difference is the difference between joining the frontier
and joining the whole closure, round after round.

\begin{keyidea}
Inference is recursive query evaluation, and semi-naive evaluation is to
recursion what a delta is to an update: only what changed can cause anything
new.
\end{keyidea}

\section{Keeping the closure current}

\subsection{Insert}

An inserted triple is a closure computation that starts small: run the same
rounds with $\Delta_0 = \{t\}$, inside the inserting transaction, so that the
closure is complete when the transaction commits. Inserting
\code{Bronze subClassOf Alloy} into a hierarchy above it derives \code{Bronze
subClassOf Metal} and \code{Bronze subClassOf Material} in two rounds, and the
types of every bronze instance in the same rounds.

\subsection{Delete is harder}

Deleting \code{Alloy subClassOf Metal} should remove \code{Alloy subClassOf
Material}, unless it still follows some other way, for example because a
curator also asserted \code{Alloy subClassOf Material} directly. A triple can
have several derivations, and deleting one support must not delete a triple
another support still holds. Two classic answers \cite{gupta1993}:

\begin{description}[leftmargin=1.2em]
\item[Counting.] Store, for each triple, how many derivations it has; a delete
decrements the counts of what it supported and removes triples whose count
reaches zero. Exact for non-recursive rules, but under recursion a triple can
support itself through a cycle, and its count never reaches zero.
\item[Delete and rederive (DRed).] Three steps: (1) \emph{overdelete}:
remove the deleted triple and, by semi-naive rounds, everything it could have
helped derive; (2) \emph{rederive}: for each overdeleted triple, check whether
it still follows in one step from what remains, and put it back if so; (3)
propagate the rederived triples by semi-naive insertion.
\end{description}

\begin{figure}[ht]
\centering
\begin{tikzpicture}[font=\scriptsize,node distance=0.55cm and 1.2cm,
  cls/.style={draw,rounded corners=2pt,minimum width=1.5cm,minimum height=0.5cm,fill=white}]
\node[cls] (al) {Alloy};
\node[cls,above=of al] (me) {Metal};
\node[cls,above=of me] (ma) {Material};
\draw[arrow,dashed,red!70!black] (al) -- node[right]{deleted} (me);
\draw[arrow] (me) -- (ma);
\draw[arrow] (al.west) to[bend left=50] node[left,align=right]{asserted:\\stays} (ma.west);
\draw[arrow,dotted] (al.east) to[bend right=50] node[right,align=left]{inferred:\\overdeleted,\\rederived} (ma.east);
\end{tikzpicture}
\caption{DRed on one delete. Overdeletion removes the inferred \code{Alloy
subClassOf Material}; rederivation finds it is also asserted, so the triple
is back before the transaction commits.}
\end{figure}

DRed is correct with recursion because rederivation asks a positive question,
``does this follow from what remains'', instead of trusting a count. Its cost is
that one delete near the top of a hierarchy overdeletes a large part of the
closure and rederives most of it.

\begin{decision}{Materialize, and maintain with DRed (planned)}
\dpart{Choice} Entailed triples are computed by semi-naive evaluation, stored
in the indexes, extended on insert, and maintained on delete by delete and
rederive. \code{\textbackslash set inference rdfs|off}, off by default.
\dpart{Alternatives} Query rewriting with property paths; counting
derivations; recomputing the closure from scratch after each delete.
\dpart{Why} Reads stay one scan and the planner stays unchanged. DRed is
correct under recursion, which counting is not, and it touches only what the
deleted triple could have supported.
\dpart{Cost} More triples and more index work per write, and a delete high in
a hierarchy is expensive. A test checks after random inserts and deletes that
the store equals a closure computed from scratch.
\end{decision}

\begin{decision}{Where inferred triples live (planned)}
\dpart{Choice} To be decided by measurement: a mark on each triple in the
existing SPO, POS, and OSP indexes, or separate permutations for inferred
triples that scans also read.
\dpart{Alternatives} The two above.
\dpart{Why} A mark keeps one scan per pattern but makes deleting all inferred
triples (\code{\textbackslash set inference off}) a full pass; separate
permutations make that a drop, but double the scans of every pattern.
\dpart{Cost} Either way, a delete must know which triples were asserted, so an
asserted triple that is also inferred keeps its asserted status.
\end{decision}

\section{What it costs}

Materialization trades space and write time for read time, and the exchange
rate depends on the hierarchy's depth. A triple deep in a hierarchy of depth
$d$ entails up to $d$ type triples, and a subclass chain of $n$ classes has
$n(n-1)/2$ entailed links. LUBM's ontology is shallow, so its closure adds a
modest fraction of triples; an anatomy ontology several levels deep adds many
more. The plan measures load time and size with and without inference on
LUBM(1), and checks that the LUBM queries that need RDFS return the official
answer counts.

\begin{status}
Planned (plan 13.1--13.5). Today OntoDB stores OWL as triples and has no
reasoner: the query for materials in \cref{ch:engine} misses \code{Alloy}, as
it should.
\end{status}

\begin{checkyourself}
\item Give one workload for which rewriting at read time beats materializing, and one for which it does not.
\item Run semi-naive evaluation on $A \sqsubseteq B$, $B \sqsubseteq C$ and the type triple \code{x a A}. Which triples are new in each round?
\item Why does the semi-naive formula join $\Delta_i$ with $T_i$ in both orders?
\item With counting, a delete in a subclass cycle $A \sqsubseteq B \sqsubseteq A$ leaves both links' counts above zero. Why, and how does DRed avoid the problem?
\item An asserted triple is also inferred. A user deletes it. What should the store hold afterwards, and what must the store know to get this right?
\end{checkyourself}

\section*{Further reading}
Abiteboul, Hull, and Vianu \cite{alicebook} on Datalog and semi-naive
evaluation; Gupta, Mumick, and Subrahmanian \cite{gupta1993} on counting and
DRed; the RDF 1.1 Semantics \cite{rdfsemantics} for the RDFS rules; Motik et
al.\ \cite{rdfox} on materialization in RDFox; Guo, Pan, and Heflin \cite{lubm}
on LUBM.
